\subsection{Monte Carlo methods}
\label{subsec:monte-carlo-methods}

Monte Carlo methods use random samples to approximate expectations.
The basic object in these notes is an expectation under some law \(p\), because once we can draw i.i.d.\ samples from \(p\), estimation becomes averaging.
Ordinary numerical integration is one example of this same idea: after choosing a sampling density, an integral can be rewritten as an expectation.
The rest of the subsection asks two practical questions: how to generate useful draws, and how to spend a fixed simulation budget well.

\subsubsection{Monte Carlo estimation}
The starting point is the law-of-large-numbers estimator for an expectation under a tractable sampling law.
A basic task is to evaluate an integral
\[
I = \int_{\mathcal{D}} g(x)\,dx,
\]
which can be rewritten as an expectation under any density $p$ on $\mathcal{D}$ that is positive wherever $g$ is nonzero:
\[
I = \E_{X \sim p}\!\left[\frac{g(X)}{p(X)}\right].
\]
In many applications we instead want
\[
\mu = \E_{X \sim p}\!\left[f(X)\right],
\]
where the target density is only known up to a normalizing constant, $p(x) = \tilde p(x)/Z$.
Thus our two recurring goals are:
(i) generate i.i.d. samples from $p$, and (ii) estimate $\mu$ from those samples.
The second task is the one we ultimately care about, but the first task is often the main computational bottleneck.
Much of Monte Carlo is about trading algorithmic effort for variance reduction so that $m$ samples are as informative as possible.

If $X_1,\dots,X_m \sim p$ i.i.d. and $\E_p[f(X)^2] < \infty$, then the Monte Carlo estimator
\[
\hat\mu_m = \frac{1}{m} \sum_{i=1}^m f(X_i)
\]
converges almost surely to $\mu$, and $\sqrt{m}(\hat\mu_m - \mu)$ is asymptotically normal with variance $\operatorname{Var}_p(f(X))$ (LLN/CLT).
This implies a standard error that scales like $1/\sqrt{m}$, so doubling accuracy typically requires roughly four times as many samples unless we use variance reduction.
Throughout this unit we assume we can numerically generate i.i.d. samples from $\mathrm{Unif}(0,1)$ and $\mathcal{N}(0,1)$.

\subsubsection{Basic sampling constructions}
\label{subsec:mc-generating}
Many Monte Carlo procedures start from the ability to sample from a target distribution.
We record two classical one-dimensional constructions that we will reuse as building blocks later.
Both methods illustrate a common theme: we trade analytic structure (closed-form CDFs or a tractable envelope) for efficient sampling.

\label{subsubsec:inverse-cdf}
The first construction pushes uniform randomness through a cumulative distribution function.
\begin{proposition}[Inverse CDF method]
Suppose $U \sim \mathrm{Unif}(0,1)$ and $F$ is a one-dimensional cumulative distribution function.
Define
\[
X = F^{-1}(U) \coloneqq \inf\{x : F(x) \ge U\}.
\]
Then $X$ has cdf $F$.
\end{proposition}
\begin{proof}
For any $x$, $\{X \le x\} = \{U \le F(x)\}$, so $\mathbb{P}(X \le x) = F(x)$.
\end{proof}
The inverse CDF method is exact and simple, but it requires either a closed form for $F^{-1}$ or a fast numerical inverse.
In higher dimensions we often rely on conditional CDFs, transformations, or acceptance-rejection instead.
See \citet[\S2.1]{liu_2002_monte_carlo} for additional discussion.

\begin{example}[Exponential distribution]
For $\mathrm{Exp}(\lambda)$, the cdf is
$F(x) = \bigl(1-e^{-\lambda x}\bigr)\,\1\{x \ge 0\}$, so
$F^{-1}(u) = -\lambda^{-1}\log(1-u)$.
If $u_i \sim \mathrm{Unif}(0,1)$, then
$x^{(i)} = -\lambda^{-1}\log(1-u_i)$ are i.i.d. $\mathrm{Exp}(\lambda)$.
\end{example}

\label{subsubsec:rejection-sampling}
The second construction samples from an easier envelope and keeps only the proposals that fall under the target.
Suppose the target density is $p(x) = \tilde p(x)/Z$ where $\tilde p$ is an unnormalized density and $Z = \int \tilde p(x)\,dx$.
Assume we can sample from a proposal density $g$ and that there is a constant $M>0$ with
$\tilde p(x) \le M g(x)$ for all $x$.
\Cref{alg:rejection-sampling} draws a proposal and accepts it with probability
$\tilde p(x)/(M g(x))$; the accepted sample has density $p$.
Indeed, the joint density of \emph{accepted} draws is proportional to $\tilde p$, and the acceptance probability is $Z/M\,.$
Equivalently, the expected number of proposals per accepted sample is $M/Z$, so choosing $g$ (and thus $M$) controls efficiency.
See \citet[\S2.2]{liu_2002_monte_carlo} for details.

\begin{algorithm}[t]
  \caption{Rejection sampling}
  \label{alg:rejection-sampling}
  \begin{algorithmic}[1]
    \Require Unnormalized target $\tilde p(x)$, proposal density $g(x)$, constant $M$ with $\tilde p(x) \le M g(x)$
    \Repeat
      \State Draw $Y \sim g$ and $U \sim \mathrm{Unif}(0,1)$.
    \Until{$U \le \tilde p(Y)/(M g(Y))$}
    \State \Return $Y$
  \end{algorithmic}
\end{algorithm}

\begin{example}[Truncated Gamma tail]
We want $X \sim \Gamma(4,1)$ conditioned on $X>10$.
An unnormalized density is
\[
\tilde p(x) = x^3 e^{-x} \1\{x>10\}.
\]
Take a shifted exponential proposal $Y = 10 + E$ with $E \sim \mathrm{Exp}(\lambda)$, i.e.
\[
 g(y) = \lambda e^{-\lambda(y-10)}\1\{y>10\}.
\]
Then
\[
\frac{\tilde p(x)}{g(x)} = \frac{x^3}{\lambda}\exp\bigl(-(1-\lambda)x - 10\lambda\bigr).
\]
Choosing $\lambda=\tfrac12$ gives
$\tilde p(x)/g(x) = 2 x^3 e^{-x/2-5}$, which is decreasing for $x\ge 10$;
thus we can take
$M = 2\cdot 10^3 e^{-10} = 2000 e^{-10}$.
The acceptance ratio for a proposed $x$ is
\[
\frac{\tilde p(x)}{M g(x)} = 10^{-3} x^3 e^{-x/2+5} \le 1,
\]
with equality at $x=10$.
\end{example}

Rejection sampling works best when the proposal closely tracks the target and $M$ is small.
In high dimensions it can become impractical because acceptance rates decay quickly, especially if $g$ underestimates the tails of $p$.

\begin{example}[A multidimensional rejection construction]
Suppose the target is the uniform distribution on the unit disk
\[
  \mathcal{D}\coloneqq \{x\in\R^2:\norm{x}_2\le 1\}.
\]
Draw \(Y\sim\mathrm{Unif}([-1,1]^2)\) and accept when \(Y\in\mathcal{D}\).
Conditioned on acceptance, the draw is uniform on \(\mathcal{D}\), and the acceptance probability is
\[
  \frac{\operatorname{area}(\mathcal{D})}{\operatorname{area}([-1,1]^2)}
  =
  \frac{\pi}{4}.
\]
This simple example is the multivariate version of the one-dimensional picture: rejection sampling is still exact, but its efficiency is the ratio between the target-set volume and the envelope volume.
For balls inside boxes that ratio deteriorates quickly with dimension, which is one concrete reason the method becomes fragile in high-dimensional problems.
\end{example}

\subsubsection{Importance sampling}
\label{subsubsec:importance-sampling}
Importance sampling estimates expectations under $p$ using samples from a more convenient proposal $g$.
Assume $\mathrm{supp}(p) \subseteq \mathrm{supp}(g)$ and define the importance weight
$w(x) = p(x)/g(x)$.
Then
\[
\mu = \E_{X \sim p}[h(X)] = \E_{X \sim g}\!\left[w(X) h(X)\right].
\]
Intuitively, $w(x)$ corrects for the fact that we are sampling from $g$ instead of $p$.
A simple estimator is
\[
\tilde\mu = \frac{1}{m}\sum_{i=1}^m w(X_i) h(X_i),
\qquad X_i \sim g.
\]
When $p(x)=\tilde p(x)/Z$ and $Z$ is unknown, we use the self-normalized estimator
\[
\hat\mu = \frac{\sum_{i=1}^m \tilde w_i h(X_i)}{\sum_{i=1}^m \tilde w_i},
\qquad \tilde w_i = \frac{\tilde p(X_i)}{g(X_i)}.
\]
The estimator $\hat\mu$ is generally biased for finite $m$, but it is consistent under standard moment conditions and is the natural choice when $Z$ is unknown.
Its finite-sample accuracy is still governed by weight variability: highly uneven weights make both the bias and the variance worse.
See \citet[\S2.5]{liu_2002_monte_carlo}.

\begin{algorithm}[t]
  \caption{Self-normalized importance sampling}
  \label{alg:importance-sampling}
  \begin{algorithmic}[1]
    \Require Unnormalized target $\tilde p(x)$, proposal density $g(x)$, test function $h(x)$
    \State Draw $X_1,\dots,X_m \sim g$ i.i.d.
    \State Compute weights $\tilde w_i = \tilde p(X_i)/g(X_i)$ and normalize
    $w_i = \tilde w_i / \sum_{j=1}^m \tilde w_j$.
    \State \Return $\hat\mu = \sum_{i=1}^m w_i h(X_i)$.
\end{algorithmic}
\end{algorithm}

A particularly useful case is when we already know how to sample from a base law \(\pi\), and then new side information \(y\) arrives.
If that information enters through a nonnegative likelihood factor \(L_y(x)\), the new target is the reweighted law
\[
  \pi^y(dx)
  \coloneqq
  \frac{L_y(x)\,\pi(dx)}{Z_y},
  \qquad
  Z_y\coloneqq \int L_y(x)\,\pi(dx).
\]
This is already interesting on its own:
it answers a concrete Monte Carlo question.
If we already have samples from \(\pi\) and then data arrive that change the target, can we recycle those samples rather than start the whole computation from scratch?

\begin{example}[Reweighting existing samples after data arrive]
\label{ex:conditioning-likelihood-reweighting}
Assume \(0<Z_y<\infty\), and let \(h\) be \(\pi^y\)-integrable.
Then
\begin{equation}
  \int h(x)\,\pi^y(dx)
  =
  \frac{\int h(x)L_y(x)\,\pi(dx)}{\int L_y(x)\,\pi(dx)}.
  \label{eq:conditioning-likelihood-reweighting}
\end{equation}
So if \(X_1,\dots,X_m\sim\pi\) are i.i.d., then
\[
  \int h(x)\,\pi^y(dx)
  \approx
  \frac{\sum_{i=1}^m L_y(X_i)\,h(X_i)}{\sum_{i=1}^m L_y(X_i)}.
\]
In other words, conditioning is a special case of self-normalized importance sampling with proposal \(\pi\), target \(\pi^y\), and unnormalized weights
\[
  \tilde w_i=L_y(X_i).
\]
The unknown constant \(Z_y\) cancels after normalization.
If \(\pi\) is a prior law and \(L_y(x)=p(y\mid x)\), then \(\pi^y\) is exactly the posterior law given \(y\), and posterior expectations can be estimated by reweighting prior samples.
\end{example}

This conditioning-by-reweighting viewpoint will reappear whenever side information changes an existing sampling problem.

Importance sampling is only as useful as the weight computations behind it.
Because those weights are ratios of potentially tiny densities, the next block collects the log-scale and transformed-density bookkeeping that will recur here, in Metropolis--Hastings acceptance ratios, in tempering, in SMC, and later in diffusion objectives.
\subsubsection{Working on the log scale}
\label{subsec:stable-evaluation}
\label{subsubsec:working-on-log-scale}

The Monte Carlo calculations above already expose the main numerical issue for sampling.
In floating-point arithmetic there are only finitely many representable numbers, so the computer does not literally manipulate the real numbers that appear in the mathematics.
Probabilities are often tiny, likelihood ratios can differ by many orders of magnitude, and a naive translation of the algebra into floating-point arithmetic can therefore overflow, underflow, or round away the smaller of two competing contributions entirely.
The log-weight checks around \Cref{alg:importance-sampling} are the first example, and the same issue will return immediately in Markov-chain acceptance ratios and transformed target densities.
Stable evaluation means rewriting the same mathematics so that the scale that matters for the algorithm remains visible to the machine.
One repeatedly tries to cancel irrelevant common factors, keep intermediate quantities in a representable range, and avoid subtractions in which all significant digits are lost.
In the sampling methods introduced so far, the algorithm usually needs a ratio, a normalization, a tail comparison, or a log normalizing constant rather than an absolute probability value.

The first recurring move is to stay on the log scale whenever the algorithm only cares about relative probabilities.
Whenever an objective or density is built from products of positive terms, the logarithm is usually the natural computational scale.
Products become sums, ratios become differences, and a common multiplicative scale becomes an additive shift that can be removed without changing any normalized quantity.
Suppose a calculation produces log weights $\ell_1,\dots,\ell_n$.
The normalized weights are then
\[
  w_i = \frac{e^{\ell_i}}{\sum_{j=1}^n e^{\ell_j}},
\]
and the same normalizing sum also determines the log normalizing constant $\log \sum_{j=1}^n e^{\ell_j}$.
Directly exponentiating the $\ell_i$ can be numerically disastrous: large positive values can overflow, while very negative values can underflow to zero before the ratios among the smaller terms are ever compared.
Writing
\[
  a \coloneqq \max_i \ell_i
\]
and subtracting this common shift before exponentiating gives
\begin{equation}
  \log \sum_{i=1}^n e^{\ell_i}
  =
  a + \log \sum_{i=1}^n e^{\ell_i-a}.
  \label{eq:log-sum-exp}
\end{equation}
This is the log-sum-exp identity.
It simply factors out the largest scale $e^a$.
After that shift, every exponent in \eqref{eq:log-sum-exp} is at most zero, so each exponential lies in $(0,1]$ and the sum is computed in a much safer floating-point range while preserving the relative sizes of the terms.

The same shift gives normalized weights directly:
\[
  w_i
  =
  \frac{e^{\ell_i-a}}{\sum_{j=1}^n e^{\ell_j-a}}.
\]
If we want an average rather than a sum, then
\[
  \log \frac{1}{n}\sum_{i=1}^n e^{\ell_i}
  =
  -\log n + \log \sum_{i=1}^n e^{\ell_i},
\]
which is often called a log-mean-exp calculation.
This pattern is what keeps self-normalized importance sampling stable when only a few draws dominate, and it is equally useful in latent-variable models where responsibilities must be renormalized at every iteration.
If only relative weights matter, one should remove the largest common scale before asking floating-point arithmetic to exponentiate anything.
In practice one usually stays on the log scale until the last step that genuinely requires ordinary probabilities.

The same issue appears when a calculation ends with a probability close to $0$ or $1$.
If a tail probability is so small that floating-point arithmetic rounds it to zero, then taking a logarithm afterwards loses all information.
For that reason one often works directly with $\log F(x)$ or $\log(1-F(x))$, depending on whether the lower or upper tail is the quantity of interest.
In statistics software these are usually exposed as log-CDF or log-survival evaluations.
Here the stable object is the logarithm of the probability, not the probability itself.
Similarly, if $F(x)$ is extremely close to $1$, then the subtraction $1-F(x)$ may round to zero before the logarithm is taken.
Stable log-survival calculations avoid that dangerous subtraction rather than trying to repair it afterwards.
More generally, subtracting two nearly equal floating-point numbers is a common source of lost precision, so algebraically equivalent forms that avoid such subtractions are usually preferable.

This matters whenever the algorithm compares tail events rather than raw densities.
Gaussian tail approximations, rejection envelopes, and likelihood-based tests can all fail in the far tails if the computation leaves the log scale too early.
The same discipline also clarifies diagnostics: a histogram of log weights is often informative precisely because the logarithm turns multiplicative instability into additive spread that we can see and reason about.
The same habit of carrying the dominant scale explicitly will reappear below when Metropolis--Hastings and Langevin-type samplers evaluate log targets and acceptance ratios.

\subsubsection{Transformed densities and Jacobians}
\label{subsubsec:transformed-densities}

Often the numerically convenient coordinates are not the scientifically meaningful ones.
We introduce an unconstrained variable \(z\), optimize or sample there, and map back to the original parameter by \(x=T(z)\).
This is useful for positive parameters, simplex coordinates, and Cholesky factors, but it changes the density: once we change coordinates, a Jacobian term must be included to keep the target correct.
The basic identity is the ordinary change-of-variables formula.
If $x=T(z)$ is a smooth bijection, then
\[
  p_Z(z) = p_X\!\bigl(T(z)\bigr)\,\lvert \det J_T(z)\rvert.
\]
Taking logarithms gives
\begin{equation}
  \log p_Z(z)
  =
  \log p_X\!\bigl(T(z)\bigr)
  +
  \log \lvert \det J_T(z)\rvert.
  \label{eq:transformed-density-log}
\end{equation}
Equation \eqref{eq:transformed-density-log} is therefore just the change-of-variables formula written on the log scale.
The transformed log density is the original log density evaluated at the mapped point, plus an additive volume correction.
This additive form is also the stable one computationally: multiplicative Jacobian factors can vary over many orders of magnitude, while on the log scale they enter as one more term in a sum.
The computational and probabilistic issues coincide here.
We change coordinates because the transformed variable $z$ may be easier to optimize over or sample in, but the Jacobian term is what makes that computational convenience represent the original model rather than a different one.

For a positive scalar parameter, the basic example is $x=e^z$.
Then \eqref{eq:transformed-density-log} becomes
\[
  \log p_Z(z) = \log p_X(e^z) + z.
\]
The exponential map keeps $x$ positive automatically, while the Jacobian term records how intervals on the $z$ scale stretch when mapped back to the original variable.
The same pattern extends to vector transforms, simplex maps, and Cholesky parameterizations.
Whenever we optimize or sample in transformed coordinates, carrying the Jacobian term on the log scale is what keeps the target mathematically correct and numerically stable.
This becomes especially important in latent-variable models and in gradient-based samplers, where transformed coordinates are often the difference between a feasible computation and an impossible one.

The same distinction matters after the computation is finished.
If an optimizer or sampler is run on an unconstrained variable $z$ but the scientific parameter of interest is $x=T(z)$, then summaries should usually be formed on the $x$ scale by transforming the retained iterates first.
For nonlinear maps, the order matters:
\[
  \E[T(Z)] \neq T(\E[Z])
\]
in general, and the same is true for credible intervals and other nonlinear summaries.
One therefore stores the draws or iterates in the computational coordinates that make the algorithm stable, but reports means, intervals, and other interpretable quantities on the model scale that the transformation defines.

Transformed objectives are also one of the most common places for implementation errors.
If gradients are coded by hand, then a useful consistency check is to compare the analytic gradient of the log density or potential at a few representative points with a finite-difference derivative of that same scalar function.
The finite-difference calculation is not meant to replace the final algorithm, and in high dimension it is far too expensive to use iteratively.
Its role is narrower and more important: it catches missing Jacobian terms, sign mistakes, and other local errors before those errors are amplified by Newton updates, Langevin proposals, or leapfrog trajectories.
These same ideas reappear whenever we sample in transformed coordinates and need the algorithmic convenience of the new parameterization to remain faithful to the original target.

With that numerical bookkeeping in place, we can ask whether the proposal is statistically well matched to the target.

\subsubsection{Proposal quality and effective sample size}
\label{subsubsec:proposal-quality-and-ess}
When $X_i \sim p$ i.i.d., a CLT yields approximate confidence intervals
$\bar h \pm 1.96\, s_h/\sqrt{m}$ for $\mu$, where
$\bar h = m^{-1}\sum_{i=1}^m h(X_i)$ and
$s_h^2 = (m-1)^{-1}\sum_{i=1}^m (h(X_i)-\bar h)^2$.
For importance sampling, proposal quality is often summarized by how concentrated the normalized weights are.
If all normalized weights are close to $1/m$, then every draw contributes comparably to the estimate.
If one or two normalized weights dominate, then only a few draws are doing most of the work.
This motivates the weight-based effective sample size proxy
\[
  \mathrm{ESS} \approx \frac{1}{\sum_{i=1}^m w_i^2},
\]
where $w_i$ now denotes the normalized importance weight.
When $p$ and $g$ are normalized so that $\E_g[w(X)]=1$, a large-sample heuristic gives
\[
  m\sum_{i=1}^m w_i^2 \approx \E_g[w(X)^2] = 1+\operatorname{Var}_g(w(X)),
\]
and therefore
\[
  \mathrm{ESS}(m) \approx \frac{m}{1+\operatorname{Var}_g(w(X))}.
\]
This is the weight-based notion of ESS: it measures how uneven the normalized importance weights are, not the autocorrelation ESS for correlated MCMC chains that appears later in \Cref{subsubsec:ess}.
(See \citet[\S2.5.3]{liu_2002_monte_carlo}.)
If $p$ and $g$ are normalized densities, then $\operatorname{Var}_g(w(X))=\chi^2(p\|g)$, where $\chi^2(p\|g)=\int (p(x)-g(x))^2/g(x)\,dx$ is the chi-squared divergence.
Thus we want $w(X)$ to be as close to constant as possible, i.e., $g(x)$ should closely match $p(x)$ in both modes and tails.
\Cref{fig:is-proposal-quality} illustrates typical outcomes.

\begin{figure}[t]
  \centering
  \begin{tikzpicture}
  \begin{axis}[
    name=plot1,
    width=0.28\textwidth,
    height=0.22\textwidth,
    xmin=-4, xmax=4,
    ymin=0, ymax=0.55,
    axis lines=none,
    ticks=none,
    domain=-4:4,
    samples=200,
  ]
    \addplot[black, thick] {0.4*exp(-x^2/2)};
    \addplot[stat221Red, thick] {0.36*exp(-x^2/(2*1.1*1.1))};
    \node[anchor=north] at (axis cs:0,0.53) {\footnotesize Good};
  \end{axis}

  \begin{axis}[
    name=plot2,
    at={(plot1.east)},
    xshift=1.1cm,
    anchor=west,
    width=0.28\textwidth,
    height=0.22\textwidth,
    xmin=-4, xmax=4,
    ymin=0, ymax=0.55,
    axis lines=none,
    ticks=none,
    domain=-4:4,
    samples=200,
  ]
    \addplot[black, thick] {0.4*exp(-x^2/2)};
    \addplot[stat221Red, thick] {0.5*exp(-(x-1.6)^2/(2*0.8*0.8))};
    \node[anchor=north] at (axis cs:0,0.53) {\footnotesize Bad};
  \end{axis}

  \begin{axis}[
    name=plot3,
    at={(plot2.east)},
    xshift=1.1cm,
    anchor=west,
    width=0.28\textwidth,
    height=0.22\textwidth,
    xmin=-4, xmax=4,
    ymin=0, ymax=0.55,
    axis lines=none,
    ticks=none,
    domain=-4:4,
    samples=200,
  ]
    \addplot[black, thick] {0.4*exp(-x^2/2)};
    \addplot[stat221Red, thick] {0.28*exp(-x^2/(2*1.6*1.6))};
    \node[anchor=north] at (axis cs:0,0.53) {\footnotesize Okay};
  \end{axis}
\end{tikzpicture}

\caption{Proposal quality in importance sampling.
  The target $p(x)$ is shown in black and the proposal $g(x)$ in red.
  Good proposals match both modes and tails, whereas poorly aligned proposals place too much mass in regions where $p(x)$ is small.
  Overly diffuse proposals are workable but less efficient.}
  \label{fig:is-proposal-quality}
\end{figure}

A useful heuristic is to choose $g$ to approximate $|h(x)|p(x)$, which controls the variance of $w(X)h(X)$.
In practice this often means matching the dominant modes and tail decay.
If $g$ is lighter-tailed than $p$, then a few extreme samples can dominate the estimator, leading to weight degeneracy.
When $p(x) \propto e^{-q(x)}$ with smooth $q$, a Laplace approximation suggests
\[
\nabla q(\mu) = 0, \qquad \Sigma = \bigl[\Hess q(\mu)\bigr]^{-1},
\]
and $g = \mathcal{N}(\mu,\Sigma)$.
For multimodal targets, a mixture proposal can capture multiple modes; for heavy tails, a $t_\nu$ proposal is often safer.

A common workflow is:
\begin{enumerate}[leftmargin=*,label=(\arabic*)]
  \item Start with a pilot proposal $g$ and a moderate sample size $m$.
  Compute a rough ESS estimate (e.g., via the sample weights) and inspect a histogram of log-weights.
  \item If weights are highly skewed, increase the variance of $g$ and/or add mixture components where large weights occur, then repeat step (1).
  \item Once the weights stabilize, run a larger batch for the final estimate.
\end{enumerate}
The histogram is often more informative than raw weights, since it makes tail behavior and outliers visible.

\subsubsection{Variance reduction}
Monte Carlo error decreases as $m$ grows, but we can often reduce variance further by exploiting structure.
The idea is to transform an unbiased estimator into another unbiased estimator with smaller variance by subtracting (or conditioning on) a related quantity whose expectation is known.
We summarize one widely used technique here; see \citet[\S2.3]{liu_2002_monte_carlo} for additional methods.
Suppose we want $\mu = \E[f(X)]$ and we have another random variable $Y$ with known expectation $\E[Y]$.
For any $b \in \R$, the control-variate estimator
\[
Z = f(X) + b\bigl(Y - \E[Y]\bigr)
\]
has the same mean as $f(X)$.
Its variance is
\[
\operatorname{Var}(Z) = \operatorname{Var}(f(X)) + b^2 \operatorname{Var}(Y) + 2b\,\operatorname{Cov}(f(X),Y).
\]
The minimizer is
\[
 b^* = -\frac{\operatorname{Cov}(f(X),Y)}{\operatorname{Var}(Y)} = -\rho_{f(X),Y}\frac{\mathrm{sd}(f(X))}{\mathrm{sd}(Y)},
\]
which yields
\[
\operatorname{Var}(Z) = \bigl(1-\rho_{f(X),Y}^2\bigr)\operatorname{Var}(f(X)) \le \operatorname{Var}(f(X)).
\]
Thus strong (negative) correlation between $f(X)$ and $Y$ produces the largest variance reduction.
This guarantee is for the population-optimal coefficient \(b^*\).
In practice one often replaces \(b^*\) by an empirical estimate \(\hat b\) computed from the same Monte Carlo sample.
That plug-in estimator is often still useful, but finite-sample variance reduction is no longer automatic: if the correlation is weak, estimating \(\hat b\) can slightly increase variance rather than decrease it.

\begin{example}[Rare-event estimation with IS + control variates]
Consider
$\theta = \E\bigl[e^X \1\{X>3\}\bigr]$ with $X \sim \mathcal{N}(0,1)$.
Naive Monte Carlo is inefficient because $\mathbb{P}(X>3)$ is tiny.
An importance-sampling alternative draws $Y \sim \mathcal{N}(\mu,1)$ for some $\mu>3$, with weight
\[
 w(y) = \frac{\phi(y)}{\phi(y-\mu)} = \exp\bigl(-\mu y + \tfrac12 \mu^2\bigr),
\]
and estimator
\[
\tilde\theta_{\mathrm{IS}} = \frac{1}{m}\sum_{i=1}^m e^{Y_i}\,\1\{Y_i>3\}\,w(Y_i),
\qquad Y_i \sim \mathcal{N}(\mu,1).
\]
Here $\phi$ denotes the standard normal density.
To further reduce variance, treat $Y$ as a control variate with $\E[Y]=\mu$.
Let
$H(y) = e^{y}\,\1\{y>3\}\,w(y)$ and define
\[
\hat\theta_{\mathrm{IS+CV}} = \frac{1}{m}\sum_{i=1}^m \Bigl(H(Y_i) + b\,(Y_i-\mu)\Bigr).
\]
A practical estimate of $b$ is
\[
\hat b = -\frac{\sum_{i=1}^m (H(Y_i)-\bar H)(Y_i-\bar Y)}{\sum_{i=1}^m (Y_i-\bar Y)^2},
\qquad \bar H = \frac{1}{m}\sum_{i=1}^m H(Y_i),\; \bar Y = \frac{1}{m}\sum_{i=1}^m Y_i.
\]
In this example, choosing $\mu \approx 4$ yields a noticeable variance reduction in practice.
Empirical estimation of $b$ can slightly increase variance, and the benefit is strongest when $H(Y)$ and $Y$ are strongly (typically negatively) correlated.
\end{example}
